% !PRINTED 565
\weberpage{565}

%<<CONT
worin die $A_{m,n}$ rationale Zahlen sind, die im Nenner nur Potenzen
von 2 und von 3 enthalten k\"onnen.

In gleicher Weise ergibt sich aus der Reihenentwickelung (6) f\"ur
die $\wp$-Funktion:

\weq{9}{\frac{g_2 g_3}{G}\, \wp(u) = \sum^{m,n} B_{m,n}\, j^{m+n} (j - 27\,.\,64)^{m+2n} \frac{w^{2\nu-2}}{(2\nu+1)!},}

worin die $B_{m,n}$ rationale Zahlkoeffizienten sind. Da\ss{} auch sie
nur Potenzen von 2 und 3 im Nenner haben, machen die ersten F\"alle
wahrscheinlich, kommt aber f\"ur uns jetzt nicht in Betracht.

Ist $g_2$ oder $g_3$ gleich Null, so ist die Substitution (7) nicht
brauchbar. Ist zun\"achst

\[
g_3 = 0,
\]

so fallen in den Reihen f\"ur $\sigma(u)$ und $\wp(u)$ die Glieder
weg, in denen $n$ positiv ist. Es ist dann

\[
% [sic] Die Quelle verwendet hier erneut a_{m,0} und b_{m,0}, laesst aber die Faktoren 2^{-m} aus (5)/(6) fort und aendert im zweiten Nenner (4m+1)! zu (2m+1)!. Die Quellennormierung bleibt unveraendert; neue A_m/B_m in (11) koennen anders normiert sein.
\sigma(u) = \sum a_{m,0}\, g_2^m\, \frac{u^{4m+1}}{(4m+1)!},
\]
\[
\wp(u) = \sum b_{m,0}\, g_2^m\, \frac{u^{4m-2}}{(2m+1)!}.
\]

Wir setzen dann

\weq{10}{u = g_2^{-\frac{1}{4}}\, w}

und erhalten

\weq{11}{\begin{aligned}
g_2^{\frac{1}{4}}\, \sigma(u) &= \sum^m A_m\, \frac{w^{4m+1}}{(4m+1)!}, \\
g_2^{-\frac{1}{2}}\, \wp(u) &= \sum^m B_m\, \frac{w^{4m-2}}{(4m+1)!},
\end{aligned}}

wenn $A_m, B_m$ wieder rationale Zahlen sind.

Ist endlich

\[
g_2 = 0,
\]

so setzen wir

\weq{12}{u = g_3^{-\frac{1}{6}}\, w,}

und erhalten

\weq{13}{\begin{aligned}
g_3^{\frac{1}{6}}\, \sigma(u) &= \sum^n A_n\, \frac{w^{6n+1}}{(6n+1)!}, \\
g_3^{-\frac{1}{3}}\, \wp(u) &= \sum^n B_n\, \frac{w^{6n-2}}{(6n+1)!}.
\end{aligned}}

In den Entwickelungen (8), (9), (11), (13) ist

\weq{14}{A_{00} = 1, \quad B_{00} = 1, \quad A_0 = 1, \quad B_0 = 1.}
